% =============================================================================
%  ch08  手术、S/T 矩阵与标架反常
%  负责人：bibliographer（PLAN_gaps G8.1--G8.5）
% =============================================================================
\section{手术、S/T 矩阵与标架反常}\label{sec:surgery}

\subsection{这一章要解决什么}

本文 \S4 的标题是“可计算性”：它要把任意闭三维流形 $M$ 上的配分函数 $Z(M)$ 化成 $S^3$
上的链环不变量。做法是两步：先证明 $Z$ 在 Dehn 手术下按一个矩阵 $K$ 变换（(4.27)、(4.28)），
再说明 $K$ 就是模群 $\mathrm{SL}(2,\ZZ)$ 在环面 Hilbert 空间上的表示，从而只需知道两个生成元
$S$ 与 $T$（\S4.4--4.5）。本文把这两步写得很紧：手术公式只用半页给出 \src{384--386}，
$S$、$T$ 矩阵只说“它们是环面微分同胚在特征标上的表示” \src{385}，
标架反常 (4.41) 只用三行推出 \src{390}。

本章要把这三块补齐，使得：
\begin{enumerate}[leftmargin=2em]
\item 你能自己把 (4.27) 从“配分函数 = 边界态的配对”推出来，并说清指标 $K_i{}^j$ 的来历（G8.1、G8.2）；
\item 你能自己验证 $S^2=(ST)^3=\mathbb 1$，并知道这个等式只在特定 $T$ 约定下精确成立（G8.3）；
\item 你能算出标架改变时 $Z$ 的相位 $e^{2\pi\ii c/24}$ 与 Wilson 线的 $e^{2\pi\ii t h_R}$，并知道两者同源（G8.4）；
\item 你能不看笔记写出 (4.51) 的推导链（G8.5），并把它与第 \ref{sec:verlinde} 节的 Verlinde 公式对上。
\end{enumerate}

\subsection{原文说了什么}

\paragraph{(i) 手术表示（p.384）。}
原文从“切开流形”讲起：沿圆圈 $C$ 的管状邻域边界 $\Sigma$ 把 $M$ 分成 $M_L$ 与 $M_R$，
两边的路径积分各给一个向量 $\psi\in\cH_L$、$\chi\in\cH_R$，$M$ 上的配分函数是它们的配对
$(\psi,\chi)$；若粘合前先用微分同胚 $K$ 作用在 $M_R$ 的边界上，配分函数变成 $(\psi,K\chi)$
\src{384}。随后原文说：“任何三维流形都可以通过对纽结作反复手术而从 $S^3$ 得到、或化归为 $S^3$
（或任何其他指定的三维流形）——这并不是一个很深刻的结果。然而，这样的描述远不唯一” \src{384}。
“远不唯一”正是 Kirby 移动的入口（G8.1）。

\paragraph{(ii) 手术公式（p.385--386）。}
在 Verlinde 基 $v_i$ 下，$K$ 由一个矩阵表示：$K\cdot v_i=\sum_jK_i{}^jv_j$ (4.25)
\src{385}；取 $M$ 上原本没有 Wilson 线（故 $\chi=v_0$ (4.24)），得
\begin{equation}
Z(\tilde M)=\sum_j K_0{}^j\,Z(M;R_j) \tag{4.27}\src{385}
\end{equation}
若手术前 $C$ 上已有 $R_i$ 线（$\chi=v_i$），则
\begin{equation}
Z(\tilde M;R_i)=\sum_j K_i{}^j\,Z(M;R_j). \tag{4.28}\src{386}
\end{equation}
原文明确说 (4.28)“稍后将在 Verlinde 猜想的一个新证明中用到” \src{386}。

\paragraph{(iii) $S$、$T$ 与环面。}
原文说 $K_i{}^j$“恰好是环面微分同胚 $K$ 在圈群不可约表示的特征标上所表示的矩阵”，
并把它归到文献 [37]（Turaev）与 [12,13]（Verlinde；Moore--Seiberg）\src{385}。
用它把 $S^2\times S^1$ 沿一条圆圈做 $S$ 手术得到 $S^3$，就得到 (4.37)、(4.38) \src{389}；
把 $S$ 换成 $T^nST^m$，得到 (4.40)；再比较两者即得标架相位 (4.41) \src{390}。

\paragraph{(iv) \S4.5 的具体计算。}
(4.33)--(4.36) 把 $S^2\times S^1$ 上的配分函数解释为 $S^2$ 上 Hilbert 空间的维数
（$\delta_{a,0}$、$g_{ab}$、$N_{abc}$）\src{388}；(4.42)--(4.44) 给出 $Z(S^3;R_j)=S_{0,j}$
与不打结 Wilson 圈的期望值 \src{391}；(4.47)--(4.51) 给出链接 Wilson 线
$L(R_i;R_j)$、$L(R_i;R_j,R_k)$ 的配分函数与 Verlinde 公式 \src{392--393}。

\paragraph{(v) 原文自己承认的边界。}
p.390 说标架系统“将在别处讨论”，并明说“在本文的简单应用中，这没有必要”：
本文所有应用只用“标准 $S$ 手术” \src{390}。这一点很重要——本章的 G8.4 是把原文略去的部分补出来，
而不是本文用到了它。

\begin{center}
\small
\begin{tabular}{@{}p{1.5cm}p{8.6cm}p{2.2cm}@{}}
\toprule
编号 & 内容 & 印刷页 \\
\midrule
(4.24)--(4.25) & $\chi=v_0$；$K\cdot v_i=\sum_jK_i{}^jv_j$ & 385 \\
(4.26)--(4.28) & 配分函数在手术下的变换；广义手术 & 385--386 \\
(4.33)--(4.36) & $S^2\times S^1$ 配分函数 = 维数：$\delta_{a,0}$、$g_{ab}$、$N_{abc}$ & 388 \\
(4.37)--(4.41) & $Z(S^3)=S_{0,0}$；$S_{mn}$；$(T^nST^m)_{0,0}$；标架相位 & 389--390 \\
(4.42)--(4.46) & $Z(S^3;R_j)=S_{0,j}$；不打结圈期望值；$S_{0,j}S_{0,k}/S_{0,0}$ & 391 \\
(4.47)--(4.51) & 链接圈的配分函数；Verlinde 公式 & 392--393 \\
\bottomrule
\end{tabular}
\end{center}

\subsection{原文没说的细节}

\paragraph{G8.1（手术表示与 Kirby 移动）}
原文说“描述远不唯一”\src{384}，但没说两个描述之间差什么。缺的是
Lickorish--Wallace 定理（任一闭可定向三维流形 = $S^3$ 沿某链环作整系数手术）与
Kirby 定理（两个手术表示给同一流形，当且仅当能由 Kirby 移动互相转换）。
更重要的：任何用手术表示定义的量，都必须证明与表示无关。本文用 $S$、$T$ 的模群关系
隐式地做到这一点，但没有把“Kirby 移动 $\leftrightarrow$ 模群关系”的对应写出来。

\paragraph{G8.2（(4.28) 的推导）}
原文只给结论与一句话理由（“$v_i$ 被替换为 $K_i{}^jv_j$”）\src{386}。
缺的是两步：为什么 $K$ 是线性表示（切割-粘合的 Hilbert 空间结构），
以及为什么 $\chi$ 恰好是 $v_0$ 或 $v_i$（实心环面无 Wilson 线时是真空态）。

\paragraph{G8.3（$S$、$T$ 的来源与关系）}
原文说 $K$ 是“环面微分同胚在特征标上的表示”，但不写 $S$、$T$ 的显式矩阵，
也不给 $S^2=(ST)^3=\mathbb 1$ 的证明（只在 p.390 一句使用）。缺的是：
$T$ 为什么对角（因为 $T:\tau\to\tau+1$ 只把每个特征标乘一个相位）、
$S$ 为什么由 Weyl--Kac 特征标公式给出、以及 $S^2=(ST)^3$ 在什么归一化下精确成立。

\paragraph{G8.4（标架反常）}
原文在 (2.25) 说大 $k$ 时的标架相位含 $c$ \src{361}，在 (4.41) 给出
$Z\to Z\exp(2\pi\ii(n-m)c/24)$ \src{390}，在 p.380 说 $t$ 重 Dehn 扭转给 Wilson 线相位
$e^{2\pi\ii t h_R}$ \src{380}。但它没说这三处的 $c$、$h_R$ 是同一个二维 CFT 的中心荷与共形权，
也没给 $h_R=(N^2-1)/(2N(N+k))$ 的来历（只在脚注 11 出现）\src{380}。

\paragraph{G8.5（(4.51) 的完整推导）}
原文的推导链是 (4.47)$\to$(4.48)$\to$(4.49)$\to$(4.50)$\to$(4.51)，其中 (4.49) 只用一句
“利用一维线性代数这一美妙事实”带过 \src{393}。缺的是把 $(\psi,\chi)(v',v)=(\psi,v)(v',\chi)$
这一步的几何来源写清楚：它是把 $L(R_i;R_j,R_k)$ 沿 $R_i$ 剪开、看成连通和。

\subsection{补出的推导}

\begin{derivation}[D8.1：手术公式 (4.27) 的四步推导]
\textbf{第 1 步（切割）。}取 $C$ 的管状邻域 $U\cong D^2\times S^1$，令 $\Sigma=\partial U$；
去掉 $U$ 的内部得 $M_L$，留下实心环 $M_R=U$。于是 $M=M_L\cup_\Sigma M_R$。

\textbf{第 2 步（配对）。}Atiyah 公理：带边流形给边界 Hilbert 空间中的向量。
$M_L$ 给 $\psi\in\cH_L$，$M_R$ 给 $\chi\in\cH_R$，而 $\cH_L\cong\cH_R^*$（因为 $\Sigma$ 在两侧定向相反），
所以 $Z(M)=(\psi,\chi)$。这一步是原文 p.384 那段话的全部内容 \src{384}。

\textbf{第 3 步（手术 = 换粘合映射）。}手术就是保留 $M_L$，把 $M_R$ 用另一个微分同胚
$K:\Sigma\to\Sigma$ 粘回去，得到 $\tilde M=M_L\cup_{K}M_R$。在 $M_R$ 一侧这等价于把态
$\chi$ 换成 $K\chi$（$K$ 作用在 $\cH_R$ 上），故 $Z(\tilde M)=(\psi,K\chi)$。

\textbf{第 4 步（用 Verlinde 基展开）。}在 $\cH_{T^2}$ 的 Verlinde 基 $v_i$ 下，
$K v_i=\sum_jK_i{}^jv_j$ (4.25)。实心环无 Wilson 线时 $\chi=v_0$ (4.24)，
故 $K\chi=\sum_jK_0{}^jv_j$，于是
\[
Z(\tilde M)=\sum_jK_0{}^j(\psi,v_j).
\]
而 $(\psi,v_j)$ 就是“在 $C$ 上放一条 $R_j$ 表示的 Wilson 线之后的 $M$ 的配分函数”
$Z(M;R_j)$——因为 $v_j$ 与 $v_0$ 的差别正是实心环中心多一条 $R_j$ 线。
代入即得 (4.27)。

\textbf{推广到 (4.28)。}若手术前 $C$ 上已有 $R_i$ 线，则 $\chi=v_i$，
$K\chi=\sum_jK_i{}^jv_j$，同样的配对给出 (4.28)。\ours
\end{derivation}

\begin{derivation}[D8.2：两种算法对比——$Z(S^3;R_j)$ 的两条路]
式 (4.43) 的推导给了我们第一个算法；\S4.2 的显式公式给了第二个。两者必须一致，
而原文用 (4.44) 做了这个核对 \src{391}。

\textbf{算法 A（先手术后求和）。}$S^3$ 由 $S^2\times S^1$ 沿圆圈 $C$ 作 $S$-手术得到。
在 $S^2\times S^1$ 上放一条与 $C$ 平行的 $R_j$ 线（图 15b），手术后它在 $S^3$ 上仍是
标准标架的不打结圈。由 (4.27)，
\[
Z(S^3;R_j)=\sum_iS_0{}^i\,Z(S^2\times S^1;R_i,R_j)
=\sum_iS_0{}^i g_{ij}=S_{0,j},
\]
中间一步用 (4.35)（两电荷给对偶度规 $g_{ij}$），最后一步用 $g$ 对称、$S$ 对称。

\textbf{算法 B（带结不变量）。}$\S4.2$ 用拆接关系直接算了 $S^3$ 上不打结圈的期望值 (4.23)：
对 $SU(N)_k$、$N$ 维表示，
\[
\langle C\rangle=\frac{q^{N/2}-q^{-N/2}}{q^{1/2}-q^{-1/2}},
\qquad q=e^{2\pi\ii/(N+k)} .
\]
与 $Z(S^3)=S_{0,0}$（(4.38)）相除，得 $Z(S^3;R)/Z(S^3)=\langle C\rangle$。
取 $N=2$、$R$ 为二维表示，$q=e^{2\pi\ii/(k+2)}$，右端 $=q^{1/2}+q^{-1/2}=2\cos\frac{\pi}{k+2}$；
而算法 A 给 $S_{0,1}/S_{0,0}=\sin\frac{2\pi}{k+2}\big/\sin\frac{\pi}{k+2}=2\cos\frac{\pi}{k+2}$。
两者相等——这就是 (4.44) \src{391}。\ours

两条路的独立性是有意义的：算法 A 只用“手术 + 边界配对 + 模矩阵”，不碰任何纽结理论；
算法 B 只用“拆接关系 + 标架”。它们给出同一个数，等于从两个方向检验了
“CS 不变量 = 纽结不变量”这个等式在最简单情形下的相容性。
\end{derivation}

\begin{derivation}[D8.3：$S$ 与 $T$ 的来源]
$T^2$ 的映射类群是 $\mathrm{SL}(2,\ZZ)$，由
$S:\tau\mapsto-1/\tau$ 与 $T:\tau\mapsto\tau+1$ 生成。亏格 $1$ 的共形块空间
$\cH_{T^2}$ 有一组基——仿射特征标 $\chi_i(\tau)$（每个可积表示 $i$ 一个），于是模群作用给出
有限阶矩阵（用 Witten 的约定，见下）：
\begin{align}
T:\ \chi_i(\tau+1)&=e^{2\pi\ii(h_i-c/24)}\chi_i(\tau)
&&\Longrightarrow&
T_i{}^j&=\delta_i{}^j e^{2\pi\ii(h_i-c/24)},\\
S:\ \chi_i(-1/\tau)&=\sum_jS_i{}^j\chi_j(\tau)
&&\Longrightarrow&
S_{ij}&\ \text{由 Weyl--Kac 特征标公式给出}.
\end{align}
$T$ 对角是定义级的：$\tau\to\tau+1$ 只把每个特征标乘一个相位，相位由 $L_0$ 在两个最高权态之间的差
$h_i-c/24$ 决定（$c$ 来自 $T(z)$ 的真空能）。$S$ 的矩阵元就是仿射特征标在模变换下的混合矩阵；
对 $SU(2)$ 它有闭式
\begin{equation}
S_{mn}=\sqrt{\frac{2}{k+2}}\,\sin\frac{(m+1)(n+1)\pi}{k+2},
\qquad m,n=0,1,\dots,k,
\label{eq:surg-s}
\end{equation}
即原文 (4.39) \src{390}（$m=2j$）。一般 $G$ 的 $S$ 由 Weyl--Kac 公式
与 $S_{\lambda\mu}=S_{\lambda 0}\,\chi_\mu\!\big(e^{-2\pi\ii(\lambda+\rho)/(k+h^\vee)}\big)$ 给出
[上野 \S5.5，书 p.103；(5.29)(5.30)]。

\emph{关系式。}在 $PSL(2,\ZZ)$ 中 $S^2=1$、$(ST)^3=1$。提升到矩阵时，
这给出 $S^2=C$（$C$ 为电荷共轭，$SU(2)$ 时 $C=\mathbb 1$）与 $(ST)^3=C$。
\textbf{关键}：这个等式是否\textbf{精确}成立，取决于 $T$ 是否含 $c/24$ 的相位。用我们的数值核对
（$SU(2)_k$，$k=1,2,3$，取 \eqref{eq:surg-s} 的 $S$）：
\begin{center}
\small
\begin{tabular}{@{}lccc@{}}
\toprule
$T$ 的取法 & $S^2$ & $(ST)^3$ & $(ST)^3$ 与 $S^2$ 的关系 \\
\midrule
$T=\mathrm{diag}(e^{2\pi\ii(h_i-c/24)})$（Witten） & $\mathbb 1$ & $\mathbb 1$ & 精确相等 \\
$T=\mathrm{diag}(e^{2\pi\ii h_i})$（拓扑自旋） & $\mathbb 1$ & $e^{2\pi\ii c/8}\mathbb 1$ & 差相位 $e^{2\pi\ii c/8}$ \\
\bottomrule
\end{tabular}
\end{center}
两种约定差一个 $e^{2\pi\ii c/8}$（因为 $T$ 差 $e^{-2\pi\ii c/24}$，三次方得 $e^{-2\pi\ii c/8}$）。
\textbf{这个相位不是数学噪声，它就是标架反常的模块化形式}：$c/24$ 是 $T$ 的相位，
$c/8$ 是 $(ST)^3$ 的相位，而 (4.41) 的 $c/24$ 是手术相位的来源——三者是同一个 $c$。\ours
\end{derivation}

\begin{derivation}[D8.4：标架反常 (4.41)]
\textbf{(a) 配分函数的相位。}用 $T^nST^m$ 代替 $S$ 从 $S^2\times S^1$ 得到 $S^3$，
则 $Z(S^3)=(T^nST^m)_{0,0}$ (4.40)。在 Verlinde 基中 $T$ 对角，
$(T^nST^m)_{0,0}=\sum_{i,j}(T^n)_{0i}S_{ij}(T^m)_{j0}$。$T$ 对角化后只剩 $i=j=0$，
故（在 $g_{00}=1$ 的归一化下）
\[
Z(S^3)=e^{2\pi\ii n(h_0-c/24)}\,S_{00}\,e^{2\pi\ii m(h_0-c/24)}
=e^{-2\pi\ii(n+m)c/24}S_{00},
\]
其中用了 $h_0=0$。与 (4.38) $Z(S^3)=S_{00}$ 相比，相位是 $e^{-2\pi\ii(n+m)c/24}$。
原文 (4.41) 写的是 $e^{2\pi\ii(n-m)c/24}$ \src{390}：差别来自右端 $T^m$ 作用在
$\cH_{T^2}$ 的\textbf{对偶}（bra）上、以及 $T$ 相位符号的约定（$q$ 与 $q^{-1}$）。
两种写法的关系是固定的：只要在 (4.40) 里一致地使用同一种约定，$|Z|$ 与 Jones 多项式的值都不变。
\textbf{本节照原文记 (4.41)，具体数值计算时按上式重算相位。}\ours

\textbf{(b) 为什么是 $c/24$。}$T$ 的相位 $e^{2\pi\ii(h_i-c/24)}$ 中的 $c/24$ 来自
$T(z)$ 的真空能：$L_0$ 在真空上的值不是 $0$ 而是 $-c/24$（在柱面上，
$T_{\text{柱}}=T_{\text{平面}}-\frac{c}{24}$）。所以任何“平移一个单位”的操作都带 $e^{-2\pi\ii c/24}$。
三维一侧的对应物是 $\eta$ 不变量/标架：改变标架 $t$ 个单位，路径积分按
$e^{2\pi\ii t c/24}$ 变换（原文 (2.25)，大 $k$ 时 $c\to d$）\src{361}。

\textbf{(c) Wilson 线的相位。}$t$ 重 Dehn 扭转作用在与表示 $R$ 相联系的边界 Hilbert 空间上，
乘以 $e^{2\pi\ii t h_R}$（原文 p.380，即 (2.33)）\src{380}：
几何上它把圈 $C$ 的标架平移 $t$ 个单位（自环绕数改变 $t$），于是“插入 $W_R(C)$ 的配分函数”
按这个相位变换。对 $SU(N)_k$、$R$ 为 $N$ 维表示，
$h_R=\dfrac{N^2-1}{2N(N+k)}$（原文脚注 11）\src{380}；这与本章 \eqref{eq:hlambda}
（第 \ref{sec:wzw} 节）给出的 $h_\lambda=\frac{(\lambda,\lambda+2\rho)}{2(k+h^\vee)}$ 是同一个公式：
取 $\lambda=\omega_1$，$(\omega_1,\omega_1)=\frac{N-1}{N}$、$2(\omega_1,\rho)=N-1$，
故 $(\lambda,\lambda+2\rho)=\frac{(N-1)(N+1)}{N}=\frac{N^2-1}{N}$，除以 $2(N+k)$ 即得。
\textbf{三处相位同源}：$c/24$（标架）、$h_R$（Wilson 线标架）、$c/8$（模关系）都来自 $L_0$ 的谱。\ours
\end{derivation}

\begin{derivation}[D8.5：(4.51) 的完整推导]
\textbf{第 1 步（两种链接）。}在 $S^2\times S^1$ 上放一条与手术圈 $C$ 平行的 $R_j$ 线，
作 $S$-手术后得到 $S^3$ 上一对\textbf{链接}圈 $L(R_i;R_j)$（$R_i$ 来自 $C$ 上原有的线）。
由 (4.28)，
\begin{equation}
Z(S^3;L(R_i;R_j))=\sum_kS_i{}^k\,Z(S^2\times S^1;R_k,R_j)=\sum_kS_i{}^kg_{kj}=S_{ij}.
\tag{4.47}
\end{equation}

\textbf{第 2 步（三条链接）。}同法，$C$ 上有 $R_i$、另两平行线上有 $R_j,R_k$：
\begin{equation}
Z(S^3;L(R_i;R_j,R_k))=\sum_mS_i{}^mN_{mjk}.
\tag{4.48}
\end{equation}

\textbf{第 3 步（独立的几何分解）。}图 17c 把 $L(R_i;R_j,R_k)$ 沿 $R_i$ 剪成两半：
它是 $L(R_i;R_j)$ 与 $L(R_i;R_k)$ 的连通和。具体地，$R_i$ 与 $R_j$ 的链接处有一对
对偶的 $R_i,\bar R_i$ 电荷，$S^2$ 上带一对对偶电荷的 Hilbert 空间是一维的，
于是可以把该处的配对按这个一维性重新组合。用到的线性代数恒等式是
\begin{equation}
(\psi,\chi)\,(v',v)=(\psi,v)\,(v',\chi),
\end{equation}
它只在 $\cH$ 一维时成立（否则是一般的矩阵乘法）。把它用在“先算 $L(R_i;R_j,R_k)$ 与
$Z(S^3;R_i)$、再与两个链接的乘积比较”上，得到
\begin{equation}
Z(S^3;L(R_i;R_j,R_k))\,Z(S^3;R_i)
=Z(S^3;L(R_i;R_j))\,Z(S^3;L(R_i;R_k)).
\tag{4.49}
\end{equation}

\textbf{第 4 步（代入）。}用 (4.47) 与 (4.43) 代入 (4.49)：
\[
Z(S^3;L(R_i;R_j,R_k))\cdot S_{0,i}=S_{ij}\,S_{ik}
\ \Longrightarrow\
Z(S^3;L(R_i;R_j,R_k))=\frac{S_{ij}S_{ik}}{S_{0,i}} .
\tag{4.50}
\]
与 (4.48) 比较即得
\begin{equation}
\frac{S_{ij}S_{ik}}{S_{0,i}}=\sum_mS_i{}^mN_{mjk}.
\tag{4.51}
\end{equation}
\textbf{读法}：把左边看成“用 $S$ 表示的量”，右边看成“用融合系数 $N_{mjk}$ 表示的量”，
(4.51) 就是“$S$ 使融合规则对角化”。若定义 $w_i=S_{0,i}\sum_mS_i{}^mv_m$，
则 (4.51) 等价于 $w_iw_j=\delta_{ij}w_j$（(4.52)--(4.54)）——幂等元。
两种读法的完整展开见第 \ref{sec:verlinde} 节，那里还会证明 (4.51) 与
$N_{ijk}=\sum_lS_{il}S_{jl}S_{kl}/S_{0l}$ 等价。\ours
\end{derivation}

\begin{derivation}[D8.6：为什么 $Z(S^2\times S^1;\cdots)$ 是维数]
(4.33)--(4.36) 是 \S4.5 的入口，但原文只写结论（“$Z(S^2\times S^1;\langle R\rangle)$ 是
$\cH_{S^2;\langle R\rangle}$ 的维数”）\src{388}。这里补出四种情形。

\textbf{一般原理。}对任何 $(2+1)$ 维 TQFT，把 $M$ 取成 $\Sigma\times S^1$，
沿 $S^1$ 方向作时间演化等价于“无演化”（$\Sigma\times S^1$ 的度规在 $S^1$ 方向平移下不变），
于是 $Z(\Sigma\times S^1)=\Tr_{\cH_\Sigma}(\mathbb 1)=\dim\cH_\Sigma$。
这就是 (4.33) 的全部内容 \src{388}。它把“配分函数”问题换成“数维数”问题，
而后者纯代数。

\textbf{四种情形。}在 $S^2$ 上放 $n$ 个标记点、分别赋表示 $R_1,\dots,R_n$，则
\begin{center}
\small
\begin{tabular}{@{}clll@{}}
\toprule
$n$ & $\cH_{S^2;R_1,\dots,R_n}$ & 维数 & 理由 \\
\midrule
0 & $\CC$ & $1$ & 无标记点的球面只有一个态（真空）\\
1 & $\Inv(R_a)$ & $\delta_{a,0}$ & 球面上单个表示的不变量只有真空表示有 \\
2 & $\Hom(R_a,R_b)$ & $\delta_{b,\bar a}$ & 等价于 $R_a\otimes\bar R_b$ 中的不变量；$g_{ab}$ 即这个配对\\
3 & $\Inv(R_a\otimes R_b\otimes R_c)$ & $N_{abc}$ & 三点共形块维数，即融合系数\\
\bottomrule
\end{tabular}
\end{center}
逐条看：$n=1$ 给 (4.34) \src{388}；$n=2$ 的对偶配对 $g_{ab}$ 给 (4.35)；
$n=3$ 给 (4.36) $Z(S^2\times S^1;R_a,R_b,R_c)=N_{abc}$ \src{388}。
四行都只用了“$\cH_{S^2}$ 的态 = $S^2$ 上表示的不变量”这一条 TQFT 公理。

\textbf{为什么这张表重要。}$g_{ab}$ 是 $\cH_{T^2}$ 上的对偶度规（用它升降指标，
把 $S_{0}{}^i$ 变成 $S_{0,j}$），$N_{abc}$ 是融合系数（用它写出 (4.36) 与 (4.51) 的右边）。
换句话说，(4.33)--(4.36) 把“几何量”$Z(S^2\times S^1;\cdots)$ 翻译成了
“表示论量”$\delta$、$g$、$N$——后面所有手术计算都靠这台翻译机。
一旦接受这张表，\S4.5 的推导就只剩 (4.47)--(4.51) 的链式代入。\ours
\end{derivation}

\begin{derivation}[D8.7：$k=1$ 的手算例子]
取 $SU(2)_1$：$c=1$，可积表示只有 $j=0,\frac12$，$h_0=0$、$h_{1/2}=\frac14$。由 \eqref{eq:surg-s}，
\[
S=\frac{1}{\sqrt2}\begin{pmatrix}1&1\\[2pt]1&-1\end{pmatrix},
\qquad
T=\begin{pmatrix}e^{-2\pi\ii/24}&0\\[2pt]0&e^{2\pi\ii(1/4-1/24)}\end{pmatrix}
=\begin{pmatrix}e^{-\ii\pi/12}&0\\[2pt]0&e^{5\ii\pi/12}\end{pmatrix}.
\]
\textbf{(a) $S^2=\mathbb 1$：}直接算 $S^2=\frac12\begin{pmatrix}1&1\\1&-1\end{pmatrix}^2=\mathbb 1$。

\textbf{(b) $(ST)^3=\mathbb 1$：}乘开得（数值核对到 $10^{-16}$）
\[
(ST)^3=\mathbb 1,
\]
但注意：若把 $T$ 换成拓扑自旋 $\mathrm{diag}(1,e^{2\pi\ii/4})=\mathrm{diag}(1,\ii)$，
则 $(ST)^3=e^{2\pi\ii/8}\mathbb 1$（因为 $c=1$，$e^{2\pi\ii c/8}=e^{\ii\pi/4}$）。
两套矩阵只差 $T\to e^{-2\pi\ii/24}T$，而三次方差 $e^{-2\pi\ii/8}$，故 $(ST)^3$ 的相位正好补偿。
这就是 D8.3 表中“两种约定差 $e^{2\pi\ii c/8}$”的具体一例。

\textbf{(c) 两个数值。}由 (4.38) 与 (4.43)，$Z(S^3)=S_{00}=\frac{1}{\sqrt2}$，
$Z(S^3;R_{1/2})=S_{0,1/2}=\frac{1}{\sqrt2}$，故不打结圈的期望值
$\langle C\rangle=Z(S^3;R_{1/2})/Z(S^3)=1$；与 (4.44) 的
$\sin\frac{2\pi}{3}\big/\sin\frac{\pi}{3}=1$ 一致。\ours

\textbf{(d) 为什么这个例子值得手算。}$k=1$ 是唯一一个 $S$ 矩阵小到可以用纸笔验证
$S^2=(ST)^3=\mathbb 1$ 的情形；一旦你在这个例子上看清 $c/24$ 的作用，
再看一般 $k$ 就不会把两种 $T$ 约定混起来。
\end{derivation}

\subsection{与其他来源的对照}

\begin{center}
\small
\begin{tabular}{@{}p{3.3cm}p{4.6cm}p{5.4cm}@{}}
\toprule
对象 & 本文 & 河野 \\
\midrule
手术不变量的定义 & (4.27)(4.28)：配分函数 = 边界态配对 & 定理 2.16（书 p.60）：$Z_k(M)=S_{00}C^{\sigma(L)}\sum_\lambda S_{0,\lambda(1)}\cdots J(L;\lambda)$ \\
Kirby 不变性 & 只说“描述远不唯一”，用模群关系隐式处理 \src{384} & 定理 2.16 的证明：Kirby 操作不变性（命题 2.15：$J(H;\lambda,\mu)=S_{\lambda\mu}/S_{00}$）\\
$S$ 矩阵来源 & 环面微分同胚在特征标上的表示（p.385） & 环面模变换（(1.42)，书 p.37 附近）\\
Verlinde 公式 & (4.51)（\S4.5） & 命题 1.22（书 p.37）\\
标架反常 & (2.25)、(4.41)（$c/24$） & 定理 2.16 中的 $C^{\sigma(L)}$（手术链环的标架/签名因子）\\
\bottomrule
\end{tabular}
\end{center}

上野 \S5.4--5.5（书 p.103--106）与本文的重叠在 $S$ 矩阵的显式公式与
$N_{ijk}=\sum_lS_{il}S_{jl}S_{kl}/S_{0l}$；Moore 的讲法（2.2.13 节，p.58）把
“曲线模空间上的物理态丛”与手术放在一起讲，适合作为第三种读法。
Schottenloher 第 11 章只讲 Verlinde 公式的几何侧，不涉及手术。

\subsection{陷阱与约定}

\begin{enumerate}[leftmargin=2em]
\item \textbf{指标位置。}$K_i{}^j$ 是线性变换的矩阵：源指标在下、像指标在上。
(4.28) 左边的 $R_i$ 与右边的求和 $j$ 是两个不同指标，不能同名缩并。
$S_{ij}$（两个下标）与 $S_i{}^j$（一上一下）数值相同只因为 $S$ 对称、且用 $g$ 升降指标；
写成 $S_{0}{}^i g_{ij}=S_{0,j}$ 时，$g$ 是把对偶指标降下来，不是 $\delta_{ij}$。
\item \textbf{$(ST)^3=S^2=1$ 是精确的还是投影的？} 用 $T=\mathrm{diag}(e^{2\pi\ii(h_i-c/24)})$
（Witten 的约定）精确成立；用拓扑自旋 $e^{2\pi\ii h_i}$ 则差 $e^{2\pi\ii c/8}$。
引用别人书里的 $S^2=(ST)^3=1$ 前，先看他用哪种 $T$。
\item \textbf{标架相位的符号。}(4.41) 写 $e^{2\pi\ii(n-m)c/24}$；直接按 (4.40) 与
$T$ 对角展开会得到 $e^{-2\pi\ii(n+m)c/24}$。差别来自右端 $T^m$ 作用在对偶上
（以及 $q$ 与 $q^{-1}$）。做具体数值时不要照抄指数，重算一遍。
\item \textbf{$Z(S^3)$ 的归一化。}$Z(S^3)=S_{0,0}$（(4.38)）不是“可以随便设成 1”的；
它被场论公理（切割-粘合相容）唯一确定。所有“期望值”都是比值
$\langle C\rangle=Z(S^3;R)/Z(S^3)$（(4.44)）。若把 $Z(S^3)$ 归一为 1，
$Z(S^3;R_j)$ 就不再等于 $S_{0,j}$，而是量子维数 $S_{0,j}/S_{0,0}$。
\item \textbf{(4.53) 的印本笔误。}原文 (4.53) 印作 $v^m$，按上下文应为 $v_m$
（本地译本已在 Tutor Note 中指出）。这不影响 (4.54)，但引用时要说明。
\item \textbf{“$S$ 手术”与“$T$ 手术”不是两种手术。}手术本身是“挖去管状邻域再粘回”，
$S$ 与 $T$ 是粘回时用的\textbf{粘合映射}（环面微分同胚）。同一个链环配不同的粘合映射，
得到的流形不同（$S$ 给 $S^3$，$T^nST^m$ 也给 $S^3$ 但标架不同）。
\end{enumerate}

\subsection{练习}

\begin{enumerate}[leftmargin=2em]
\item \textbf{（复算）} 用 \eqref{eq:surg-s} 写出 $k=2$ 的 $3\times3$ $S$ 矩阵，
取 $T=\mathrm{diag}(e^{2\pi\ii(h_i-c/24)})$、$c=3/2$、$h=(0,\frac{3}{16},\frac12)$，
验证 $S^2=\mathbb 1$ 与 $(ST)^3=\mathbb 1$。
\item \textbf{（复算）} 用 (4.43) 与 (4.39) 算 $k=1$ 时二维表示 Wilson 圈的期望值
$\langle C\rangle$，并说明为什么它等于 $1$（提示：$k=1$ 时只有一个非平凡可积表示，
$S_{0,1}/S_{0,0}=2\cos\frac{\pi}{3}=1$）。
\item \textbf{（推导）} 从 (4.40) 与 $T$ 对角出发重推标架相位，说明为什么
原文写 $n-m$ 而我们展开得到 $-(n+m)$；给出两种约定之间的换算。
\item \textbf{（推导）} 由 (4.47) 与 (4.43) 推出 (4.50)，再由 (4.48) 推出 (4.51)，
并指出哪一步用了 $g_{kj}$ 的对偶性。
\item \textbf{（推导）} 证明 (4.49)：把 $L(R_i;R_j,R_k)$ 沿 $R_i$ 剪开，
用 $\cH$ 一维时的恒等式 $(\psi,\chi)(v',v)=(\psi,v)(v',\chi)$；说明为什么 $\cH$ 必须一维。
\item \textbf{（判断）} “(4.28) 就是 Verlinde 公式。”——对还是错？请说明 (4.28) 与 (4.51)
各自把 $S$ 用在哪里。
\item \textbf{（判断）} “标架反常只影响整体相位，所以可以忽略。”——对还是错？
（提示：$|Z|$ 不变，但 Wilson 线的链接数、手术表示的选择都会改变相位；
Jones 多项式作为 $q$ 的多项式对相位敏感，因为相位里含 $q$ 的幂。）
\end{enumerate}

\subsection{本章边界}

\begin{boundary}
\begin{itemize}[leftmargin=1.4em]
\item \textbf{原文给出}：(4.27)(4.28) 及推导思路（p.384--386）、$K$ 是模群表示（p.385）、
(4.37)--(4.41)（p.389--390）、(4.42)--(4.51)（p.391--393）、标架修正 (2.25)（p.361）与
Wilson 线标架 (2.33)（p.380）。
\item \textbf{我们补的}：D8.1 的第 1--4 步显式化；D8.2 的两种算法对比；
$S$、$T$ 矩阵的显式构造与 $S^2=(ST)^3$ 的两种约定数值核对；
(4.41) 的符号分析；$h_R$ 与 $h_\lambda$ 的同一性；D8.5 的五步推导与一维 Hilbert 空间的用法。
\item \textbf{未证}：Kirby 定理与 Lickorish--Wallace 定理（只引用，未给证明）；
$S$ 矩阵的 Weyl--Kac 推导（只给结果与 $SU(2)$ 闭式，一般 $G$ 未推）；
一般手术下 $K=T^nST^m$ 的完整分类（只对 $S$ 与 $T^nST^m$ 两种情形给出）；
(4.51) 与 $\sum_lS_{il}S_{jl}S_{kl}/S_{0l}$ 的等价性留到第 \ref{sec:verlinde} 节。
\item \textbf{数值核对的范围}：$S^2=(ST)^3$ 的表只对 $SU(2)$、$k=1,2,3$ 核过；
一般 $G$ 的 $S^2=C$（$C\neq\mathbb 1$）未数值核对。
\end{itemize}
\end{boundary}
